

\documentclass[lighthipster]{simplehipstercv}

\usepackage[utf8]{inputenc}
\usepackage[default]{raleway}
\usepackage[margin=1cm, a4paper]{geometry}
%------------------------------------------------------------------ Variablen

\newlength{\rightcolwidth}
\newlength{\leftcolwidth}
\setlength{\leftcolwidth}{0.23\textwidth}
\setlength{\rightcolwidth}{0.75\textwidth}


\title{New Simple CV}
\author{\LaTeX{} Ninja}
\date{June 2019}

\pagestyle{empty}
\begin{document}
\documentclass[lighthipster]{simplehipstercv}
\usepackage[utf8]{inputenc}
\usepackage[default]{raleway}
\usepackage[margin=1.5cm, a4paper]{geometry}
\usepackage{array}

%------------------------------------------------------------------ Variables
\newlength{\rightcolwidth}
\newlength{\leftcolwidth}
\setlength{\leftcolwidth}{0.23\textwidth}
\setlength{\rightcolwidth}{0.75\textwidth}

%------------------------------------------------------------------
\title{CV Emanuel Gómez}
\author{Emanuel}
\date{Marzo 2026}

\pagestyle{empty}
\begin{document}

\thispagestyle{empty}


\simpleheader{headercolour}{\textcolor{white}{Emanuel}}{\textcolor{white}{Gómez}}{}{}

\subsection*{}
\vspace{2em} 

\setlength{\columnsep}{1cm}
\columnratio{0.25}[0.70]
\begin{paracol}{2}
\hbadness5000

\footnotesize
{\setasidefontcolour
\begin{center}
    \roundpic{foto.jpeg} 
\end{center}

\bigskip\bigskip

\begin{flushleft}
    \infobubble{\faEnvelope}{cvgreen}{white}{emaesteban2003@gmail.com}
    \infobubble{\faPhone}{cvgreen}{white}{+54 261 5584339} \\[0.8em]
    \infobubble{\faGithub}{cvgreen}{white}{EmanuelGomezRastrilla} \\[0.8em]
    
\end{flushleft}
}

\switchcolumn

\vspace{2em} 
\small
\section*{Curriculum Vitae}
\justifying
{Mi nombre es Emanuel Gómez Rastrilla, nací el 17 de Julio del 2003, en Mendoza, Argentina donde actualmente sigo viviendo. Considero que tengo una facilidad para socializar y integrar a todos a los grupos de trabajo aunque no me destaco por el rol de líder.}

\vspace{2.5em}

\section*{Idiomas}
\vspace{0.5em}
Empecé en un instituto de inglés alrededor de los 10 años , luego de algunos años en el intituto ECI de inglés me cambié el instituto Mother Teresa de Maipú, en donde hice 4 años más pero me ví obligado a dejar por mi cursado de la escuela secundaria la cual era doble turno, más el pre universitario. Tengo muchos años practicando el inglés aunque nunca pude conseguir un título y a día de hoy sigo teniendo un buen nivel.

\vspace{2.5em}


\section*{Estudios}
\vspace{0.5em}
\begin{tabular}{@{} p{2.8cm} p{8cm} @{}}
    \bfseries 2009--2022 & 
    \textbf{Escuela Técnica Industrial Emilio Civit (Maipú)} \newline
    \textsl{Técnico Electricista} \newline
    Realicé mis estudios secundarios en la Escuela Técnica Industrial Emilio Civit (Maipú), donde obtuve el título de Técnico Electricista. Durante esta etapa adquirí una sólida formación en electricidad, electrónica básica, instalaciones eléctricas y fundamentos técnicos aplicados al ámbito industrial.\\
    
    \multicolumn{2}{c}{\vspace{1em}} \\ 

    \bfseries 2023--Actualidad & 
    Actualmente me encuentro cursando Ingeniería Industrial en la Universidad Nacional de Cuyo, donde estoy desarrollando mi formación en áreas como gestión de procesos, optimización de sistemas productivos, análisis de datos y organización industrial.\\
\end{tabular}

\vspace{2.5em}

\section*{Proyectos y Experiencia}
\vspace{0.5em}
\textbf{Instalación Eléctrica Residencial – Proyecto Práctico} \newline
\textbf

 Diseño y montaje de una instalación eléctrica domiciliaria, incluyendo distribución de circuitos, instalación de tableros, protecciones eléctricas y cableado conforme a normas básicas de seguridad eléctrica. 

\end{paracol}
\end{document}

